\clearpage
% ============================================================================
\section*{v4.10.0-r70 Source-Only: Port Weights against Amplitudes; the Compensator on the State Sphere; Correction to r69.4 (Exploration k382)}
\addcontentsline{toc}{section}{v4.10.0-r70 Source-Only: Port Weights against Amplitudes; the Compensator on the State Sphere; Correction to r69.4 (Exploration k382)}
\label{sec:v4100-r70-port-weights}
% ============================================================================

\paragraph{Purpose.}
Remark v4.10.0-r69.4 gave the matter frame the port weights \(1:x\) (\(76.8\,\%\) in the compensator port) and, in the same sentence, the
ellipse of the explorations k292 and k297 with \(b/a=x\).  A drawing of the frames requested by the author on 26 November 2026 showed that
these are two different states: \(b/a=x\) has the weights \(1:x^2\) (\(91.65\,\%\)).  This layer states where the frames lie on the sphere of
field states, what the compensator does there, which of the two statements of r69.4 stands, and how the numbers of the notes change that
inserted \(x\) as a ratio of amplitudes.  \textbf{Exploration, not blind}: every statement is an identity or a recalculation; there is no test
with a decision criterion and no H~number.  No status of an open problem changes and no guard changes.  Earlier layers are not rewritten.

\noindent\textbf{Source-only status.}
This v4.10.0-r70 layer is source-only.  It does not build a PDF, does not assign a new DOI,
and does not authorize a public release.  The v5.5.0 Version DOI
\href{https://doi.org/10.5281/zenodo.24563811}{10.5281/zenodo.24563811} and the Concept DOI
\href{https://doi.org/10.5281/zenodo.19210728}{10.5281/zenodo.19210728} are unchanged by it.

\subsection*{v4.10.0-r70 the state sphere of the doublet}
\label{subsec:v4100-r70-sphere}

Notation of v4.9.0-r04: \(F=[\mathbf E\mid\mathbf B]\), \(G=F^{\mathsf T}F\), \(u=\tfrac12\operatorname{tr}G\),
\(\rho=G/\operatorname{tr}G=\tfrac12(1+b_S\sigma_z+b_C\sigma_x)\) with
\[
b_S=\frac{|\mathbf E|^2-|\mathbf B|^2}{2u},\qquad b_C=\frac{\mathbf E\cdot\mathbf B}{u},\qquad b_3:=\frac{|\mathbf E\times\mathbf B|}{u}\ \ge0 .
\]
\(b_S\) and \(b_C\) are \(b_1\) and \(b_2\) of r04.2.  The compensator of the doublet is \(C_{\HC}=\tfrac12(1+\sigma_x)\) (r04.1), so its axis is
\(b_C\) (``compensator axis along \(x\)'', v4.9.0-r27).

\paragraph{Theoremlet v4.10.0-r70.1 (hemisphere of states and position of the frames).}
\begin{enumerate}[leftmargin=2em]
\item[(a)] \(b_S^2+b_C^2+b_3^2=1\).  Every field state is a point of the hemisphere \(b_3\ge0\) of the unit sphere, and the state disk of
r04.2 is its projection along \(b_3\).  The third coordinate is the share of the energy that flows.
\item[(b)] A boost of rapidity \(\eta\) transverse to the field takes the state \(\mathbf E=\mathbf B\) (\(b_C=1\)) to
\(b_C=\operatorname{sech}2\eta\), \(b_S=0\), \(b_3=\tanh2\eta\): it moves along the meridian from the pole \(b_C=1\) towards the point
\(b_3=1\), which it reaches only for \(\eta\to\infty\).
\item[(c)] With the port weights of Convention~F (``compensator port holding \(p_+\)'') the frames are R at the pole \(b_C=1\)
(\(\mathbf E=\mathbf B\)), M at \(b_C=1/\gL\), that is \(\eta=\varepsilon=\operatorname{artanh}\bsym\) and the angle
\(\arccos(1/\gL)=57.57^\circ\) from the pole, P at \(b_3=1\) (\(\mathbf E\perp\mathbf B\), \(|\mathbf E|=|\mathbf B|\)), and T at \(b_C=-1/\gL\).
In the representative with \(|\mathbf E|=|\mathbf B|\) the angle between \(\mathbf E\) and \(\mathbf B\) equals the angle from the pole:
\(0^\circ\), \(57.57^\circ\), \(90^\circ\), \(122.43^\circ\).
\item[(d)] On this meridian \(D\) (\(b\mapsto-b\)) and \(\vartheta\) (\(b_C\mapsto-b_C\)) both exchange M and T, and \(J_B\)
(\(b_S\mapsto-b_S\)) fixes all four frames.
\end{enumerate}

\paragraph{Proof sketch.}
(a) \(|b|^2=1-\det G/u^2\) and \(\det G=|\mathbf E\times\mathbf B|^2\) (r04.2).  (b) \(\mathbf E'=\gamma(\mathbf E+\mathbf v\times\mathbf B)\),
\(\mathbf B'=\gamma(\mathbf B-\mathbf v\times\mathbf E)\) with \(\mathbf v\perp\mathbf E=\mathbf B\); \(\mathbf E\cdot\mathbf B\) is invariant and
\(u\mapsto u\cosh2\eta\).  (c) \(p_\pm=(1\pm b_C)/2\) and \(p_+=1/(1+x)\) give \(b_C=(1-x)/(1+x)=1/\gL=\operatorname{sech}2\varepsilon\)
(v4.9.0-r24.3, v4.10.0-r66).  (d) v4.9.0-r04.2(b), r28.  Checked exactly with sympy and independently in Wolfram Language (k382b, W1--W3).
\(\square\)

\subsection*{v4.10.0-r70 the compensator on the sphere}
\label{subsec:v4100-r70-compensator}

\paragraph{Theoremlet v4.10.0-r70.2 (the compensator is the measurement along the axis \(b_C\)).}
\begin{enumerate}[leftmargin=2em]
\item[(a)] The share of the energy in the compensator port is \(p_+=(1+b_C)/2=|\mathbf E+\mathbf B|^2/(4u)\), that of the partner port
\(p_-=(1-b_C)/2\).  The compensator sees of a state only its height \(b_C\).
\item[(b)] The kept part \(FC_{\HC}\) has \(\mathbf E'=\mathbf B'=(\mathbf E+\mathbf B)/2\) and lies at the pole \(b_C=1\), the frame R; the
removed part lies at the opposite pole \(b_C=-1\) (v4.10.0-r23.2 for the kept part).
\item[(c)] The non-selective step \(\rho\mapsto C_{\HC}\rho C_{\HC}+A\rho A\) of v4.9.0-r27 is the projection \((b_S,b_C)\mapsto(0,b_C)\)
onto the axis.
\item[(d)] Two routes lead from R to the height of M.  The boost of rapidity \(\varepsilon\) (r70.1(b)) reaches the point M on the sphere.
The duality rotation by \(\alpha\) with \(\tan\alpha=\bsym\), which moves R along the rim \(b_3=0\) by \(2\alpha\), followed by the
non-selective step, reaches \(b_C=\cos2\alpha=1/\gL\) (r27.1(b)); here \(p_+=\cos^2\alpha=0.7681\).  The compensator does not tell the two
apart.  They differ in \(b_3\): \(\tanh2\varepsilon=0.8440\) on the first route, and no single field state on the second.
\end{enumerate}

\paragraph{Proof sketch.}
(a) \(\operatorname{tr}(GC_{\HC})=u+\mathbf E\cdot\mathbf B\) (r04.3).  (b), (c) direct.  (d) \(\cos(2\arctan\beta)=(1-\beta^2)/(1+\beta^2)\).
Sympy and Wolfram (k382b, W4--W6).  \(\square\)

\paragraph{Reading.}
``A frame is the distribution of one state over the two ports'' (r69.4) means: a frame is a height on the axis of the compensator.  The
sector rule of v4.8.0-r07 (matter \(\operatorname{Re}s>0\), that is \(b_S>0\)) uses the other axis of the disk.  The two axes are related by a
duality rotation by \(45^\circ\) (v4.8.0-r07.2); which axis Convention~F means is fixed only by the words ``compensator port''.  This layer
follows these words and r27.  \textbf{Open; not decided here.}

\subsection*{v4.10.0-r70 weights against amplitudes}
\label{subsec:v4100-r70-weights}

\paragraph{Theoremlet v4.10.0-r70.3 (the two roles of \(x\)).}
Let \(a\,e^{i\omega t}+b\,e^{-i\omega t}\) be the oscillation of k253 and k292 with real \(a>b\ge0\), weights \(a^2:b^2\), Beltrami coefficient
\(\mu=b/a\), axis ratio \(K=(1+\mu)/(1-\mu)\) and height \(h=(a^2-b^2)/(a^2+b^2)\) on the sphere of k292.
\begin{enumerate}[leftmargin=2em]
\item[(a)] Weights \(1:x\) (Convention~F): \(\mu=\bsym\), \(K=e^{2\varepsilon}=3.4386\), \(h=1/\gL\), angle \(57.57^\circ\).  The height is that of
M in r70.1(c).
\item[(b)] Amplitudes \(1:x\) (k292, k297): weights \(1:x^2\), that is \(91.65\,\%:8.35\,\%\), \(K=\gL\), \(h=(1-x^2)/(1+x^2)=0.8330\), angle
\(33.59^\circ\).  This is the height after a boost with velocity \(x=\bsym^2\), not with \(\bsym\).
\item[(c)] The two agree only for \(x\in\{0,1\}\), that is at R and at P.
\end{enumerate}

\paragraph{Proof sketch.}
\(h=(1-\mu^2)/(1+\mu^2)\); \((1+\beta)/(1-\beta)=e^{2\operatorname{artanh}\beta}\); \(\operatorname{sech}(2\operatorname{artanh}x)=(1-x^2)/(1+x^2)\).
Sympy and Wolfram (k382b, W7, W8).  \(\square\)

\paragraph{Correction v4.10.0-r70.4 (to Remark v4.10.0-r69.4 and to the status of Theorem~\ref{thm:hc-quasiconformal}).}
\begin{enumerate}[leftmargin=2em]
\item[(a)] In r69.4 the sentence on the port weights stands.  The sentence that the matter frame is the ellipse with \(b/a=\mu=x\) and axis
ratio \(\gL\) does not describe the same state.  If the two counter-rotating parts are read as the ports (k253), the matter frame of
Convention~F has \(b/a=\bsym\) and the axis ratio \(e^{2\varepsilon}=3.4386\); \(\gL=\cosh2\varepsilon\) is then the inverse of its height
\(b_C\) and not its axis ratio.  The readings ``circle at R'', ``degenerate ellipse at P'' and ``reversed sense of rotation at T'' are not
affected.
\item[(b)] The proof of Theorem~\ref{thm:hc-quasiconformal} calls \(\mu=\bsym^2\) ``the ratio of the eigenvalues'', a ratio of weights, and
inserts it into \(K=(1+|\mu|)/(1-|\mu|)\), which holds for a ratio of amplitudes.  The identity \((1+x)/(1-x)=\gL\) is correct.  The statement
``dilatation \(\gL\)'' holds for the state of r70.3(b) and not for the matter frame of Convention~F.  A status note is added to the theorem.
Theorem~\ref{thm:tribonacci-beltrami}, which rests on the same \(\mu\), inherits the remark.
\item[(c)] The explorations k297 and k350 call the exchange \(a\leftrightarrow b\) (\(\mu\mapsto1/\mu\)) the exchange \(J\).  With the two parts
read as the ports, an exchange of the ports is \(D\) or \(\vartheta\) in Convention~F as corrected in v4.9.0-r28, and the involution of the
compensator fixes both parts (r70.1(d)).  Result~1 of H159 (``the image of \(C\) is the set \(|a|=|b|\)'') therefore holds for the projector
built from the exchange \(a\leftrightarrow b\), and not for the compensator of the doublet, whose image is R (r70.2(b)).  The typing of the
exchange in the picture of k253 is \textbf{open}.
\end{enumerate}

\subsection*{v4.10.0-r70 recalculation}
\label{subsec:v4100-r70-recalculation}

\paragraph{Theoremlet v4.10.0-r70.5 (exploration k382, not blind).}
Nine scripts of the notes of 24 to 26 November 2026 insert \(x\) for \(b/a\): k292, k297, k298 (H132), k299, k301, k302 (H133), k303, k305
(H136), k350 (H159).  Each was run unchanged and with \(b/a=\bsym=0.5494056\); the files were not altered.
\begin{center}
\small
\begin{tabular}{@{}>{\raggedright\arraybackslash}p{0.44\textwidth}>{\raggedright\arraybackslash}p{0.22\textwidth}>{\raggedright\arraybackslash}p{0.24\textwidth}@{}}
\toprule
quantity & \(b/a=x\) (as recorded) & \(b/a=\bsym\) (weights of Convention~F)\\
\midrule
share of the smaller part in the norm (k297) & \(8.35\,\%\) & \(23.19\,\%\)\\
angle on the sphere of k292 & \(33.6^\circ\) & \(57.57^\circ\)\\
axis ratio \(K\) of the ellipse & \(\gL=1.8647\) & \(e^{2\varepsilon}=3.4386\)\\
\(\lambda_B/\lambda_A=\mu^{-4}\) (k297) & \(120.5\) & \(10.98\)\\
\(m_H^2/(\eta f^2)=8\mu^2\) (k297) & \(0.729\) & \(2.415\)\\
\(m_s^2/m_H^2=\eta\mu^2/\pi^2\) (k297) & \(0.00923\,\eta\) & \(0.03058\,\eta\)\\
\(m_-/m_+=\mu\) for \(g_L=g_R\) (H132) & \(0.3018\) & \(0.5494\)\\
\(\rho\) in version~S with \(1+\xi=1+\mu^2\) (k299) & \(0.9165\) & \(0.7681\)\\
\(T=1/\cos^22\theta\) against the window \([5/3,3]\) (k301) & \(1.441\), below & \(3.477\), above\\
\(\kappa=1/\cos2\theta\) against the target \(1.2039\) (k303) & \(1.2004\) (\(-0.29\,\%\)) & \(\gL=1.8647\) (\(+55\,\%\))\\
\bottomrule
\end{tabular}
\end{center}
\begin{enumerate}[leftmargin=2em]
\item[(a)] No verdict changes.  H132 (no massless gauge field on the compensator direction), H133, H136 and H159 rest on structure that does
not depend on the value of \(\mu\); in H159 none of the ratios becomes \(\gL^5\) (exponents between \(-1.98\) and \(+2.98\)).
\item[(b)] With \(b/a=\bsym\) the angle of k292 equals the angle of M in r70.1(c) and the angle \(2\alpha\) of v4.10.0-r66.  The number that k292
recorded as ``computed only, no assignment'' is, in this reading, the position of the matter frame.
\item[(c)] The only numerical hit of the series, \(\kappa=(\gL+1/\gL)/2\) (k303; with the reference value of \(\bsym\) it is \(1.20049\),
\(-0.28\,\%\) from the target), exists only for \(b/a=x\).  It is \(1/\cos2\theta\) with \(\tan\theta=\mu\); for \(b/a=\bsym\) the same expression is
\(\gL\), and no candidate of the list of k303 lies within \(1\,\%\) of the target (nearest: \(1/\sin2\theta=1.1848\), \(-1.6\,\%\)).  Whether the hit
is chance, as k303 already allowed, or an indication that the state of r70.3(b) has a meaning of its own, is not decided.
\item[(d)] In k350 three ratios become exact powers: norm R\,:\,M at equal area \(=\gL^{-1}\), charge R\,:\,M at equal \(f\) \(=\gL\), and
\(K_1\) R\,:\,M at equal \(f\) \(=\gL^{2}\).  They restate \(h=1/\gL\).  The near agreement \(\cos^22\theta\approx1-\bsym^2\) (\(0.6940\) against \(0.6982\))
recorded there disappears (\(0.2877\) against \(0.4507\)).
\end{enumerate}

\paragraph{Check.}
\texttt{k382\_nachrechnung.py} runs the nine scripts twice and prints the changed lines (\texttt{k382\_out.txt}); the closed forms were
recomputed independently in Wolfram Language (k382b, W9--W12).  The recorded output of k297 has \(m_H^2=0.7285\) and \(\lambda_B/\lambda_A=120.6\)
because it used \(\gL=1.8644\); with the reference value these are \(0.7289\) and \(120.5\).

\subsection*{v4.10.0-r70 frames}
\label{subsec:v4100-r70-frames}

\begin{center}
\small
\begin{tabular}{@{}>{\raggedright\arraybackslash}p{0.12\textwidth}>{\raggedright\arraybackslash}p{0.14\textwidth}>{\raggedright\arraybackslash}p{0.12\textwidth}>{\raggedright\arraybackslash}p{0.16\textwidth}>{\raggedright\arraybackslash}p{0.30\textwidth}@{}}
\toprule
frame & \(b_C\) & \(b_3\) & \(p_+:p_-\) & angle between \(\mathbf E\) and \(\mathbf B\) (\(|\mathbf E|=|\mathbf B|\))\\
\midrule
R & \(1\) & \(0\) & \(1:0\) & \(0^\circ\)\\
M & \(0.5363\) & \(0.8440\) & \(0.7681:0.2319\) & \(57.57^\circ\)\\
P & \(0\) & \(1\) & \(\tfrac12:\tfrac12\) & \(90^\circ\)\\
T & \(-0.5363\) & \(0.8440\) & \(0.2319:0.7681\) & \(122.43^\circ\)\\
\bottomrule
\end{tabular}
\end{center}
No new frame map is used.  r70.1 and r70.2 are frame statements of Convention~F in the names of v4.10.0-r69; r70.3 to r70.5 concern the matter
frame and, through \(\mu\mapsto1/\mu\), the dual frame.  The values depend on \(\bsym\), which is inserted (reference value of v4.10.0-r67).

\subsection*{v4.10.0-r70 register entries}
\label{subsec:v4100-r70-register}

\begin{center}
\small
\begin{tabular}{@{}>{\raggedright\arraybackslash}p{0.40\textwidth}>{\raggedright\arraybackslash}p{0.52\textwidth}@{}}
\toprule
statement & status after v4.10.0-r70\\
\midrule
frames as points of the hemisphere \(b_S^2+b_C^2+b_3^2=1\) (r70.1) & exact; positions depend on the inserted \(\bsym\)\\
compensator as measurement along \(b_C\) (r70.2) & exact; restates v4.9.0-r04.3, r27 and v4.10.0-r23.2 on the sphere\\
r69.4, ellipse of the matter frame with \(b/a=x\), axis ratio \(\gL\) & corrected: that state has the weights \(1:x^2\); with the weights of Convention~F \(b/a=\bsym\), axes \(e^{2\varepsilon}:1\) (r70.4(a))\\
Theorem~\ref{thm:hc-quasiconformal}, dilatation \(\gL\) & status note: holds for a ratio of amplitudes \(x\), not for the matter frame of Convention~F (r70.4(b))\\
numbers of k292, k297, k298, k299, k301, k302, k303, k305, k350 & typed as values at \(b/a=x\); values at \(b/a=\bsym\) in r70.5; no verdict changes\\
number hit \(\kappa=(\gL+1/\gL)/2\) (k303) & belongs to \(b/a=x\) only; chance or indication not decided (r70.5(c))\\
which axis of the disk carries the frames; typing of the exchange \(a\leftrightarrow b\) & open (reading after r70.2; r70.4(c))\\
OP1--OP5 & open, as before\\
\bottomrule
\end{tabular}
\end{center}

\subsection*{v4.10.0-r70 notes}
\label{subsec:v4100-r70-notes}

{\raggedright
Notes: \texttt{notes/\allowbreak research/\allowbreak v4\_\allowbreak 10\_\allowbreak 0\_\allowbreak r70\_\allowbreak portgewichte\_\allowbreak hc/}
(k382, k382b with outputs) and \texttt{notes/\allowbreak research/\allowbreak frames\_\allowbreak kugelbild\_\allowbreak 2026\_\allowbreak 11\_\allowbreak 26/}
(drawing in four panels with its check).
Tools: Python with sympy (exact) for the identities and for running the nine scripts; Wolfram Language independently for the identities and
the closed forms.  Qiskit not used (no circuit).  No HDL/MMA comparison (no \texttt{RUN\_TAPE} sequence).  The Wolfram counter-checks of the nine
scripts themselves (k292b and the others) were not rerun.
\par}

\subsection*{v4.10.0-r70 ledger and guard update}
\label{subsec:v4100-r70-ledger-guard}

\[
R_{\rm chk}^{v4100\text{-}r70}=(1,1,0)
\]
for (status recorded; computation checked in two tools; no comparison with data).

The layer spends no guard:
\[
R_{\rm glob}^{(v4100\text{-}r70,+)}=(0,1,1,1,1),
\qquad
\Gamma_{\rm guard}^{v4100\text{-}r70}=(1,1,1,1,1),
\qquad
\beta_4=1/256\ \text{unspent}.
\]
The layer does not imply exact finite \(C5\) zero closure, exact all-readout terminality, \(C6/T2\), a global photon-clock/RH theorem,
or global all-readout closure.  OP1--OP5, C5 and C6/T2 remain open.
